\chapter{Considerações Finais}
\label{chap:conclusao}

Esta dissertação investigou o Prêmio de Risco de Volatilidade do Bitcoin (BVRP), definido como a diferença entre a volatilidade realizada em janela retrospectiva de 30 dias e a volatilidade implícita extraída do índice DVOL da Deribit, sobre uma amostra diária de 1.775 observações entre março de 2021 e março de 2026. O objetivo foi caracterizar o prêmio em três dimensões complementares: sua existência e magnitude, seu conteúdo preditivo sobre retornos futuros e sua utilidade econômica em estratégias condicionais. As seções seguintes retomam as quatro hipóteses formuladas no Capítulo~\ref{chap:introducao}, explicitam as limitações da análise e apontam extensões.

\section{Síntese das evidências}
\label{sec:conclusao_sintese}

\textbf{O prêmio existe e é economicamente relevante (H1).} O BVRP médio da amostra é de aproximadamente $-8{,}25$ pontos percentuais, com a volatilidade implícita situando-se sistematicamente acima da realizada. A magnitude e o sinal são consistentes com a interpretação clássica de prêmio de seguro pago pelos compradores de opções, replicando em criptoativos um padrão amplamente documentado em mercados de \emph{equity}. Este é o achado mais robusto da dissertação: não depende de especificação econométrica, mantém-se em todos os subperíodos anuais analisados e é insensível ao sinal do retorno corrente.

\textbf{O prêmio não prediz linearmente os retornos futuros (H2).} Regressões em seis horizontes (1 a 60 dias) produzem coeficientes uniformemente não significativos, com o menor valor-$p$ atingindo 0,241 --- resultado robusto à reamostragem em blocos e à correção para múltiplos testes. A ausência de previsibilidade linear é, portanto, uma conclusão negativa firme, e não um artefato de baixa potência em uma única especificação. O Capítulo~\ref{chap:regimes} reforça esse achado ao mostrar que nem mesmo a interação entre o prêmio e o regime de volatilidade recupera significância estatística convencional.

\textbf{A utilidade econômica existe, mas é defensiva (H3).} Estratégias condicionais a quantis extremos do BVRP reduzem o \emph{drawdown} máximo de aproximadamente $-78\%$ para $-38\%$, operando em apenas 22,5\% do tempo no limiar de referência. O ganho, contudo, não se traduz em Sharpe superior ao do compra e manutenção no mesmo período operacional: o sinal opera como filtro de risco, excluindo estados de mercado adversos, e não como preditor de direção. O padrão é robusto ao longo dos quantis testados (60\% a 90\%) e sobrevive a custos de transação de até 10 \emph{bps} por operação, ainda que a proteção de cauda se deteriore modestamente nesse cenário.

\textbf{A não linearidade não resgata a previsibilidade, mas a dimensão condicional é inequívoca (H4).} Modelos não paramétricos alcançam $R^2$ fora da amostra de 0,42 a 0,44 na previsão do próprio BVRP --- praticamente indistinguíveis dos modelos lineares e muito aquém do \emph{benchmark} de persistência ($R^2 = 0{,}94$), que nenhum modelo supervisionado supera. Cabe registrar que essa comparação é estruturalmente assimétrica: o \emph{benchmark} utiliza o valor contínuo do prêmio corrente, enquanto os modelos o acessam apenas de forma discretizada, em três faixas. A assimetria, contudo, é reveladora --- é justamente essa variável discretizada que concentra quase toda a capacidade preditiva dos modelos, enquanto retornos defasados, variações de volatilidade e calendário contribuem de forma marginal. Em outras palavras, a informação sobre o prêmio de amanhã está essencialmente no prêmio de hoje, e a categorização em faixas destrói a maior parte dela. A dimensão condicional, por sua vez, é confirmada pela análise do Capítulo~\ref{chap:regimes}: o prêmio médio varia de $-13{,}20$ p.p. em regimes de baixa volatilidade a $+0{,}66$ p.p. em regimes de alta, praticamente desaparecendo sob estresse.

\textbf{Leitura conjunta.} O conteúdo informacional do BVRP não reside na previsão pontual de seu nível --- inatingível além da própria persistência da série --- nem na antecipação da direção dos retornos, mas na identificação do estado de volatilidade do mercado. O prêmio é grande, persistente e sistematicamente negativo; é essa regularidade, e não sua variação de curto prazo, que carrega significado econômico. A dissipação do prêmio em regimes de estresse, documentada no Capítulo~\ref{chap:regimes}, é a manifestação mais nítida desse mecanismo: o seguro é caro justamente quando o sinistro não ocorre.

\section{Limitações}
\label{sec:conclusao_limitacoes}

\textbf{Extensão amostral.} A amostra cobre cinco anos, período curto para inferência sobre prêmios de risco em um ativo de volatilidade estrutural elevada. A baixa precisão estatística das estimativas --- visível nos intervalos de confiança amplos do Capítulo~\ref{chap:resultados_ols} --- é consequência direta dessa limitação, e não deve ser confundida com evidência de ausência de efeito.

\textbf{Quebra estrutural associada aos ETFs \emph{spot}.} A aprovação dos ETFs \emph{spot} de Bitcoin pela SEC em janeiro de 2024 ocorreu dentro do período amostral e introduziu canal institucional regulado de exposição ao ativo, com potencial efeito sobre a demanda por proteção via opções. A análise trata a amostra como bloco único, hipótese que se justifica pela transição gradual do mercado, mas que não foi testada formalmente.

\textbf{Pulverização amostral na análise por regimes.} A decomposição das estratégias condicionais por regime de volatilidade realizada reduz drasticamente o número de observações com sinal ativo em cada célula --- em particular no regime de baixa volatilidade, com apenas 23 observações. Os resultados dessa decomposição são reportados como descritivos e não sustentam inferência confiável.

\textbf{Escolhas de calibração.} O quantil de entrada das estratégias é selecionado a partir do conjunto de limiares testados sobre a amostra completa; embora o limiar seja recalculado em janela expansiva a cada instante, a escolha do percentil de referência utiliza informação de todo o período. Analogamente, os hiperparâmetros dos modelos baseados em árvores adotam valores convencionais, sem otimização sistemática --- pendência reconhecida no Capítulo~\ref{chap:cap5_ml}.

\textbf{Fricções não modeladas.} Os resultados incorporam custos de transação unilaterais de até 10 \emph{bps} por operação, mas não capturam \emph{slippage} em períodos de baixa profundidade, restrições de posicionamento ou fricções de liquidez intradiária, que podem reduzir os ganhos documentados em implementações reais.

\section{Extensões para pesquisas futuras}
\label{sec:conclusao_extensoes}

\textbf{Componentes de memória longa.} Os testes de estacionariedade do Capítulo~\ref{chap:dados} indicam que RV, IV e BVRP exibem padrão compatível com processos altamente persistentes --- ADF rejeita a raiz unitária, mas KPSS também rejeita a estacionariedade. A incorporação de componentes no espírito do modelo HAR (médias mensais e trimestrais) como preditores adicionais é a extensão mais natural, e permitiria avaliar se a persistência de longo prazo carrega informação além das defasagens curtas empregadas nesta dissertação.

\textbf{Seleção de variáveis em dois estágios.} A combinação de um estágio de seleção via LASSO com um modelo não paramétrico no segundo estágio permitiria separar o efeito da seleção de preditores do efeito da flexibilidade funcional, hoje confundidos na comparação entre modelos lineares regularizados e modelos baseados em árvores.

\textbf{Teste formal da quebra estrutural.} A hipótese de estabilidade dos parâmetros ao longo do período amostral pode ser testada por procedimentos como Chow, CUSUM ou janelas móveis de regressão centradas na data de aprovação dos ETFs. A estimação separada para os subperíodos pré e pós-ETF permitiria avaliar se o sinal e a magnitude do prêmio de seguro se modificaram após o ingresso institucional pleno.

\textbf{Previsão de magnitude.} A evidência preliminar reportada no apêndice do Capítulo~\ref{chap:regimes} --- de que o BVRP se relaciona com a magnitude, e não com a direção, do movimento futuro --- é economicamente coerente com a interpretação de prêmio de seguro, mas estatisticamente marginal sob a correção mais conservadora para retornos sobrepostos. Amostras mais longas permitiriam avaliar se o sinal se consolida.

\textbf{Estrutura a termo do prêmio.} Esta dissertação restringe-se ao horizonte de 30 dias, alinhado ao índice DVOL. A extensão para múltiplos vencimentos permitiria caracterizar a estrutura a termo do prêmio e verificar se a dissipação em regimes de estresse é uniforme ao longo da curva ou concentrada em vencimentos curtos.

\section{Observação final}
\label{sec:conclusao_final}

O resultado central desta dissertação é, em larga medida, negativo: o BVRP do Bitcoin não prevê retornos. Essa conclusão tem valor próprio. A literatura de criptoativos convive com um viés de publicação favorável a achados positivos, frequentemente obtidos sob protocolos que não controlam viés de antecipação, sobreposição de janelas ou múltiplos testes. Documentar com rigor que um sinal economicamente plausível \emph{não} carrega previsibilidade direcional --- e simultaneamente identificar onde seu conteúdo informacional de fato reside --- é contribuição tão legítima quanto a demonstração de um efeito. O prêmio de risco de volatilidade do Bitcoin é real, é grande e é informativo sobre o estado do mercado; não é, contudo, um preditor de retornos.
